
\renewcommand{\Folder}{2025_MapaMunicipiosTamaulipas}
\renewcommand{\EntradaBibtex}{2025_MapaMunicipiosTamaulipas}

\begin{frame}{\citetitle{\EntradaBibtex} \footnotemark[1] (1)}
%\begin{block}{Brazo Robótico \footnotemark} 
\begin{columns}
	\begin{column}{0.58\textwidth}
		\begin{itemize}
			\item A partir de una API de Mapas y de Polígonos de municipios del INEGI, se implementó una de MAPAS
			\item El usuario navega sobre el mapa de tamaulipas y selecciona un punto de su interes
		\end{itemize}
	\end{column}
	\begin{column}{0.2\textwidth}
		 \begin{center}
		     \includegraphics[width=0.98\textwidth]{\Folder/figs/initialScreen}
		 \end{center}
	\end{column}
	\begin{column}{0.2\textwidth}
		 \begin{center}
		     \includegraphics[width=0.98\textwidth]{\Folder/figs/clickingMunicipality}
		 \end{center}

	%\end{column}
	%\begin{column}{0.22\textwidth}  
		%\begin{center}
		 %\includegraphics[width=0.98\textwidth]{\Folder/figs/cuadroG}
		 %\end{center}
	\end{column}
\end{columns}
%\end{block} 
\footnotetext[1]{\fullcite{\EntradaBibtex}}
\end{frame}


\begin{frame}{\citetitle{\EntradaBibtex} \footnotemark[1] (2)}
%\begin{block}{Brazo Robótico \footnotemark} 
\begin{columns}
	\begin{column}{0.58\textwidth}
		\begin{itemize}
			\item Se resaltan las atracciones específicas de ese
			\item Se despliega información del municipio seleccionado
		\end{itemize}
	\end{column}
	\begin{column}{0.2\textwidth}
		 \begin{center}
		     \includegraphics[width=0.98\textwidth]{\Folder/figs/openMunicipalityPOI}
		 \end{center}
	\end{column}
	\begin{column}{0.2\textwidth}
		 \begin{center}
		     \includegraphics[width=0.98\textwidth]{\Folder/figs/shownMunicipalityDetails}
		 \end{center}

	%\end{column}
	%\begin{column}{0.22\textwidth}  
		%\begin{center}
		 %\includegraphics[width=0.98\textwidth]{\Folder/figs/cuadroG}
		 %\end{center}
	\end{column}
\end{columns}
%\end{block} 
%\footnotetext[1]{\fullcite{\EntradaBibtex}}
\end{frame}

